% ============================================================
% PARTE VII - TONICIZZAZIONE E MODULAZIONE
% CAPITOLO 33 - TONICIZZAZIONE E DOMINANTI SECONDARIE
% ============================================================
\part{Tonicizzazione e modulazione}

\chapter{Tonicizzazione e dominanti secondarie}
\label{chap:tonicizzazione}

La tonicizzazione è il primo passo verso la modulazione: un grado non tonicale viene trattato per un momento come se fosse la tonica.

\section{Modulazione e tonicizzazione}
\label{sec:mod_vs_ton}

\begin{definizione}[Modulazione]
    È il cambio \textbf{stabile} di centro tonale, confermato da una cadenza nella nuova tonalità (Capitolo \ref{chap:modulazione}).
\end{definizione}

\begin{definizione}[Tonicizzazione]
    È il trattamento \textbf{momentaneo} di un grado non tonicale come tonica locale, senza abbandonare la tonalità principale. Il ritorno alla tonica originale avviene subito dopo, entro uno o due accordi.
\end{definizione}

\begin{regola}
    Se l'accordo tonicizzato è seguito da una cadenza nella nuova tonalità locale si ha una modulazione. Se è seguito da un ritorno alla tonalità originale si ha una tonicizzazione.
\end{regola}

\section{La dominante secondaria}
\label{sec:meccanismo_dominante}

\begin{definizione}[Dominante secondaria]
    È un accordo costruito e funzionante come dominante di un grado diverso dalla tonica. Si indica come V/II («dominante del secondo»), V/V («dominante del quinto») e così via. È sempre una triade maggiore o una settima di dominante.
\end{definizione}

\begin{regola}
    Si può tonicizzare qualsiasi grado tranne il VII (e, nel minore, il II), perché la triade diminuita non può fare da tonica. I casi più comuni sono il II, il IV e il V.
\end{regola}

Due casi sono particolarmente frequenti.

\section{V/V: dominante della dominante}
\begin{itemize}
    \item Struttura: accordo maggiore sul II grado, con la terza alzata (Re-Fa$\sharp$-La in Do maggiore).
    \item Risoluzione: su V o \sigla{V}{7}{}.
    \item Condotta: il Fa$\sharp$, sensibile della dominante, sale a Sol.
\end{itemize}

\begin{esempio}[V/V in Do maggiore]
	\label{ex:v_di_v}
	\musica{es_cap15_01}
\end{esempio}

\section{V/IV: dominante della sottodominante}
\begin{itemize}
    \item Struttura: accordo sul I grado con settima minore (Do con \Sibemolle{} in Do maggiore).
    \item Risoluzione: su IV.
\end{itemize}

\begin{esempio}[V/IV in Do maggiore]
	\label{ex:v_di_iv}
	\musica{es_cap15_02}
\end{esempio}

\section{Riconoscimento nel basso dato}
\label{sec:riconoscimento}

\begin{regola}
    Se nel basso compare un'alterazione non prevista dall'armatura ci si chiede: questa nota è la sensibile (la terza) di una dominante secondaria? I segnali tipici sono il Fa$\sharp$ al basso (IV grado alzato, terza del V/V in primo rivolto) e il \Sibemolle{} al basso (VII grado abbassato, settima del V/IV in terzo rivolto).
\end{regola}
